% homework/hw05x/hw05x.tex
% Year 7 - Extra Homework: two-digit multiplication, and long division.
%
% Goes out alongside Homework 5. Same framing rule as hw03x and hw04x: nothing
% on the sheet says easy, basic, extra help, catch-up or revision, and no
% question is marked as being for anybody in particular. The cover states a
% fact, and it is a true one -- every method in Unit 1 ends in a multiplication
% or a division. Factors are found by dividing. The HCF is a division. A square
% past 12 x 12 is a two-digit multiplication and a cube is two of them. So the
% two operations on this sheet are not beside Unit 1; they are what Unit 1 is
% made of.
%
% This is the sequel to hw04x and picks up exactly where it stopped:
%   * hw04x taught two rows, then add, and divided by a two-digit number by
%     counting up in multiples. This sheet teaches the WRITTEN long division --
%     divide, multiply, subtract, bring down -- which is the method that does
%     not fall over on a four-digit dividend.
%   * The multiplication here goes one step further too: three digits by two,
%     which is what a cube is.
%
% It is a drill sheet on purpose. Long stretches of the same operation, working
% boxes big enough to set the sum out properly, and a check -- multiply back --
% attached to every division section.
%
% House style is shared: see homework/hw-style.tex. Method: the new-homework
% skill. Source is pure ASCII on purpose.
%
% Build:  pwsh .claude/skills/new-homework/scripts/build-hw.ps1 hw05x

\documentclass[12pt,a4paper]{article}

% -- Answer key switch -------------------------------------------------------
\newif\ifkey
\ifdefined\TEACHER\keytrue\else
  \keyfalse        % \keyfalse = student copy   |   \keytrue = teacher copy
\fi

% -- House style -------------------------------------------------------------
% homework/hw-style.tex
% Shared house style for every Year 7 homework packet. \input this from the
% packet file, right after \documentclass and the answer-key switch.
%
% Do not fork this file per packet. If a packet needs a different look, the
% question is whether every packet should have it.
%
% The choices here are deliberate and were arrived at by printing and looking:
%   * No solid dark banners. Every header is a light tint with a coloured spine
%     on the left, so colour reads clearly without flooding a school printer.
%   * No marks and no mark totals. These are practice, not exam papers.
%   * Lato at 12pt with open leading, plus mathastext so inline maths is Lato
%     too. Without mathastext every "-6 + 4" is Computer Modern serif and the
%     sheet looks half-typeset.
%   * Writing lines are spaced for Year 7 handwriting, not adult margin notes.
%
% Provides: \wlines \blank \tickbox \sep \qmk-free marks (none) \hwsection
% \qhead \expr, the workspace box, and the remember / wordhelp / activity /
% yourturn coloured boxes. Palette matches the lesson decks exactly.

% -- Packages ----------------------------------------------------------------
\usepackage[T1]{fontenc}
\usepackage[utf8]{inputenc}
\usepackage[default]{lato}
\usepackage[a4paper,top=1.7cm,bottom=1.7cm,left=2.0cm,right=2.0cm,headheight=15pt]{geometry}
\usepackage{amsmath,amssymb}
\usepackage{xcolor}
\usepackage[most]{tcolorbox}
\usepackage{enumitem}
\usepackage{tikz}
\usetikzlibrary{arrows.meta}
\usepackage{array}
\usepackage{tabularx}
\usepackage{fancyhdr}
\usepackage{multicol}
\usepackage{needspace}
\usepackage{graphicx}

% Make maths use Lato too. Without this, every "-6 + 4" on the page is set in
% Computer Modern serif and the sheet looks half-typeset. Must come last.
\usepackage[italic]{mathastext}

\linespread{1.06}\selectfont

% -- The Learner's Book palette (same hex values as the lesson decks) ---------
\definecolor{cbteal}{HTML}{0087A8}
\definecolor{cbpurple}{HTML}{5C2483}
\definecolor{cborange}{HTML}{C25E12}
\definecolor{cbgreen}{HTML}{4A8B23}
\definecolor{cbred}{HTML}{C8102E}
\definecolor{cbblue}{HTML}{1A5FA8}
\definecolor{rulegray}{HTML}{9AA5AD}

% -- Answer lines and blanks -------------------------------------------------
% \wlines{n} draws n ruled writing lines, spaced wide enough for Year 7
% handwriting rather than for an adult with a fine pen.
\makeatletter
\newcommand{\wlines}[1]{%
  \par\vspace{0.75em}%
  \@tempcnta=0\relax
  \@whilenum\@tempcnta<#1\do{%
    \hbox to \linewidth{\color{rulegray}\leaders\hrule height 0.5pt\hfill}%
    \vspace{1.5em}%
    \advance\@tempcnta by 1\relax}%
  \par\vspace{-0.8em}%
}
\makeatother

% \blank[width] is an inline gap to write a short answer in.
\newcommand{\blank}[1][2.6cm]{\underline{\hspace{#1}}}

% A boxed space for showing working.
\newtcolorbox{workspace}[1][3.4cm]{%
  colback=white, colframe=rulegray, boxrule=0.5pt, arc=5pt,
  height=#1, left=6pt, right=6pt, top=4pt, bottom=4pt,
  before skip=9pt, after skip=11pt}

% An empty tick box.
\newcommand{\tickbox}{\fbox{\rule{0pt}{1.15em}\rule{1.15em}{0pt}}}

% A middot separator.
\newcommand{\sep}{\,\textperiodcentered\,}

% -- Section and question headings -------------------------------------------
% Light tint plus a fat coloured spine on the left. No solid dark banner: the
% colour reads clearly but barely costs the printer anything.
% needspace stops a heading stranding alone at the foot of a page. Kept modest:
% too large and it pushes headings over, leaving a worse hole than it prevents.
\newcommand{\hwsection}[2][cbteal]{%
  \needspace{8\baselineskip}%
  \par\vspace{13pt}%
  \noindent
  \begin{tcolorbox}[enhanced, nobeforeafter, width=\textwidth,
      colback=#1!7, colframe=#1!7, arc=5pt,
      borderline west={5pt}{0pt}{#1},
      left=14pt, right=10pt, top=7pt, bottom=7pt]
    {\large\bfseries\color{#1}#2}
  \end{tcolorbox}%
  \par\vspace{9pt}}

% \qhead[colour]{A2}{Moving along the line} -- a question heading.
\newcommand{\qhead}[3][cbteal]{%
  \needspace{6\baselineskip}%
  \par\vspace{11pt}%
  \noindent{\large\bfseries\color{#1}#2}\quad{\large\bfseries #3}\par\vspace{4pt}}

% -- Coloured boxes ----------------------------------------------------------
% Shared look: rounded, light fill, coloured rule, coloured title text on a
% pale title strip rather than a dark one.
\tcbset{hwbox/.style={
  enhanced, breakable, arc=6pt, boxrule=1pt,
  left=11pt, right=11pt, top=8pt, bottom=8pt,
  fonttitle=\bfseries, attach boxed title to top left={xshift=12pt, yshift=-8pt},
  boxed title style={arc=4pt, boxrule=1pt, left=7pt, right=7pt, top=2pt, bottom=2pt},
  top=13pt}}

% Orange: things to remember. Matches the deck's "Write This Down".
\newtcolorbox{remember}[1][Remember from the lesson]{hwbox,
  colback=cborange!6, colframe=cborange, coltitle=cborange,
  colbacktitle=white, title={#1}}

% Blue: English help. The barrier in this class is the wording, not the maths.
\newtcolorbox{wordhelp}[1][Word help]{hwbox,
  colback=cbblue!6, colframe=cbblue, coltitle=cbblue,
  colbacktitle=white, title={#1}}

% Purple: an activity or instruction.
\newtcolorbox{activity}[1][How to do this homework]{hwbox,
  colback=cbpurple!6, colframe=cbpurple, coltitle=cbpurple,
  colbacktitle=white, title={#1}}

% Green: your turn / a friendly nudge.
\newtcolorbox{yourturn}[1][Your turn]{hwbox,
  colback=cbgreen!7, colframe=cbgreen, coltitle=cbgreen!75!black,
  colbacktitle=white, title={#1}}

% -- Shorthands --------------------------------------------------------------
\newcommand{\dC}{\,^{\circ}\mathrm{C}}          % use inside maths: $-8\dC$
\newcommand{\key}[1]{{\color{cborange}\textbf{#1}}}
\newcommand{\eng}[1]{{\color{cbblue}\textbf{#1}}}

\setlist[enumerate]{itemsep=8pt, topsep=5pt, leftmargin=*}
\setlist[itemize]{itemsep=4pt, topsep=4pt, leftmargin=*}
\renewcommand{\arraystretch}{1.45}

% -- An expression the student is asked to work from -------------------------
\newcommand{\expr}[1]{%
  \tcbox[on line, colback=cbgreen!10, colframe=cbgreen!55, boxrule=1pt, arc=4pt,
         left=9pt, right=9pt, top=4pt, bottom=4pt]{\large\bfseries $#1$}}

% -- Running head and foot ---------------------------------------------------
% The packet overrides \fancyhead[L] with its own title.
\pagestyle{fancy}
\fancyhf{}
\fancyhead[L]{\small\color{black!55}Year 7\sep Homework}
\fancyhead[R]{\small\color{black!55}Name: \blank[3.4cm]}
\fancyfoot[C]{\small\color{black!55}Page \thepage}
\renewcommand{\headrulewidth}{0pt}
\renewcommand{\headrule}{\vspace{2pt}\hbox to\headwidth{\color{cbteal!45}\leaders\hrule height 1.2pt\hfill}}



% -- Per-packet running head -------------------------------------------------
\fancyhead[L]{\small\color{black!55}Year 7\sep Extra Homework}

\begin{document}

% ===========================================================================
% COVER BLOCK
% ===========================================================================
\thispagestyle{fancy}

\begin{tcolorbox}[enhanced, colback=cbgreen!8, colframe=cbgreen!8, arc=7pt,
                  borderline west={6pt}{0pt}{cbgreen},
                  left=18pt, right=14pt, top=11pt, bottom=11pt]
  {\color{cbgreen!75!black}\bfseries YEAR 7\sep UNIT 1\sep EXTRA HOMEWORK}\\[6pt]
  {\color{cbgreen!70!black}\Huge\bfseries Bring Down the Next Digit}\\[10pt]
  {\color{black!70}Two-digit multiplication\sep and long division, written out}
\end{tcolorbox}

\vspace{7pt}
\noindent
\begin{tabularx}{\textwidth}{@{}X X@{}}
  \textbf{Name:} \blank[5.6cm] & \textbf{Class:} \blank[5.6cm] \\
  \textbf{Date given:} \blank[4.6cm] & \textbf{Date due:} \blank[5.0cm] \\
\end{tabularx}

\vspace{5pt}
\begin{activity}[Why this sheet exists]
  Every method in Unit 1 finishes in one of \key{two} operations. You find a
  factor by \key{dividing}. You find the HCF by dividing twice. You check a
  square past $12 \times 12$ by \key{multiplying}, and a cube by multiplying
  twice.

  \medskip
  Last time you divided by counting up in multiples. That works until the number
  gets big. This sheet is the \key{written method} --- the one that does not fall
  over on 4 digits.

  \medskip
  \eng{No calculator}, and show your working. Working you can see is working you
  can check.
\end{activity}

% ===========================================================================
% PART 1 -- THE FACTS
% ===========================================================================
\hwsection[cbgreen]{Part 1 \textperiodcentered\ The facts both methods run on}

\qhead[cbgreen]{E1}{Fill in the missing number}

\noindent Every long multiplication and every long division is these facts, one
after another. Aim to know them without counting.

\vspace{7pt}
\begin{multicols}{4}
\begin{enumerate}[label=(\alph*), itemsep=9pt]
  \item $7 \times 6 = \blank[1.1cm]$
  \item $8 \times 7 = \blank[1.1cm]$
  \item $9 \times 6 = \blank[1.1cm]$
  \item $6 \times 8 = \blank[1.1cm]$
  \item $9 \times 8 = \blank[1.1cm]$
  \item $7 \times 7 = \blank[1.1cm]$
  \item $4 \times 9 = \blank[1.1cm]$
  \item $8 \times 8 = \blank[1.1cm]$
  \item $9 \times 7 = \blank[1.1cm]$
  \item $6 \times 6 = \blank[1.1cm]$
  \item $9 \times 9 = \blank[1.1cm]$
  \item $8 \times 4 = \blank[1.1cm]$
  \item $12 \times 7 = \blank[1.1cm]$
  \item $11 \times 8 = \blank[1.1cm]$
  \item $12 \times 9 = \blank[1.1cm]$
  \item $7 \times 4 = \blank[1.1cm]$
\end{enumerate}
\end{multicols}

% -- E2 ---------------------------------------------------------------------
\qhead[cbgreen]{E2}{Multiply by a ten}

\begin{remember}[The second row is always a ten]
  To multiply by \key{20}: multiply by \key{2}, then \key{put a zero on the end}.
  $34 \times 2 = 68$, so $34 \times 20 = 680$. That is the whole trick, and it is
  what the second row of a long multiplication is.
\end{remember}

\vspace{2pt}
\begin{multicols}{3}
\begin{enumerate}[label=(\alph*), itemsep=9pt]
  \item $26 \times 10 = \blank[1.6cm]$
  \item $26 \times 30 = \blank[1.6cm]$
  \item $47 \times 20 = \blank[1.6cm]$
  \item $52 \times 40 = \blank[1.6cm]$
  \item $18 \times 60 = \blank[1.6cm]$
  \item $73 \times 50 = \blank[1.6cm]$
\end{enumerate}
\end{multicols}

% -- E3 ---------------------------------------------------------------------
\qhead[cbgreen]{E3}{Count up in a two-digit number}

\begin{remember}[Do this before you divide]
  Before dividing by 23, write out the first few multiples of 23. Then the
  division is only \key{reading your own list}. This is two minutes that saves
  ten.
\end{remember}

\vspace{4pt}
\noindent
\begin{tabularx}{\textwidth}{|p{1.5cm}|X|X|X|X|X|X|}
  \hline
  \textbf{$\times$} & \textbf{1} & \textbf{2} & \textbf{3} & \textbf{4} & \textbf{5} & \textbf{6} \\
  \hline
  $23$ \rule[-0.7em]{0pt}{2.2em} & & & & & & \\ \hline
  $16$ \rule[-0.7em]{0pt}{2.2em} & & & & & & \\ \hline
  $24$ \rule[-0.7em]{0pt}{2.2em} & & & & & & \\ \hline
  $37$ \rule[-0.7em]{0pt}{2.2em} & & & & & & \\ \hline
\end{tabularx}

% ===========================================================================
% PART 2 -- MULTIPLICATION
% ===========================================================================
\hwsection[cbgreen]{Part 2 \textperiodcentered\ Two rows, then add}

\begin{remember}[The method, once, in full]
  \noindent
  \begin{minipage}[t]{0.28\linewidth}
    \vspace{2pt}
    {\renewcommand{\arraystretch}{1.05}%
    \begin{tabular}{r@{\ }r@{\ }r@{\ }r}
              &   & 5 & 8 \\
      $\times$&   & 3 & 4 \\[1pt]
      \hline
              & 2 & 3 & 2 \\
      1       & 7 & 4 & 0 \\[1pt]
      \hline
      1       & 9 & 7 & 2 \\[1pt]
      \hline
    \end{tabular}}
  \end{minipage}%
  \begin{minipage}[t]{0.72\linewidth}
    \vspace{2pt}
    \key{Row 1:} $58 \times 4 = 232$. \\
    \key{Row 2:} $58 \times 30 = 1740$ --- multiply by 3, put the zero on. \\
    \key{Add them:} $232 + 1740 = 1972$. \\[4pt]
    \eng{One row means you forgot the tens.} Your answer comes out about ten times
    too small. \textbf{Two rows, every time.}
  \end{minipage}
\end{remember}

% -- E4 ---------------------------------------------------------------------
\qhead[cbgreen]{E4}{Warm up: two digits by one digit}

\vspace{2pt}
\noindent
\begin{tabularx}{\textwidth}{|X|X|X|X|}
  \hline
  $53 \times 7$ & $68 \times 6$ & $87 \times 9$ & $76 \times 8$ \\
  \rule[-2.6em]{0pt}{3.2em} & & & \\
  \hline
\end{tabularx}

\vspace{9pt}
\noindent
\begin{tabularx}{\textwidth}{|X|X|X|X|}
  \hline
  $94 \times 4$ & $47 \times 8$ & $65 \times 7$ & $89 \times 6$ \\
  \rule[-2.6em]{0pt}{3.2em} & & & \\
  \hline
\end{tabularx}

\vspace{9pt}
\noindent
\begin{tabularx}{\textwidth}{|X|X|X|X|}
  \hline
  $58 \times 9$ & $73 \times 6$ & $96 \times 7$ & $84 \times 5$ \\
  \rule[-2.6em]{0pt}{3.2em} & & & \\
  \hline
\end{tabularx}

% -- E5 ---------------------------------------------------------------------
\qhead[cbgreen]{E5}{Two digits by two digits}

\noindent \key{Two rows, then add.} Show both rows every time.

\vspace{4pt}
\noindent
\begin{tabularx}{\textwidth}{|X|X|X|}
  \hline
  $32 \times 14$ & $45 \times 23$ & $28 \times 36$ \\
  \rule[-3.3em]{0pt}{4.1em} & & \\
  \hline
\end{tabularx}

\vspace{9pt}
\noindent
\begin{tabularx}{\textwidth}{|X|X|X|}
  \hline
  $57 \times 42$ & $63 \times 27$ & $86 \times 35$ \\
  \rule[-3.3em]{0pt}{4.1em} & & \\
  \hline
\end{tabularx}

\vspace{9pt}
\noindent
\begin{tabularx}{\textwidth}{|X|X|X|}
  \hline
  $74 \times 58$ & $49 \times 61$ & $95 \times 24$ \\
  \rule[-3.3em]{0pt}{4.1em} & & \\
  \hline
\end{tabularx}

\vspace{9pt}
\noindent
\begin{tabularx}{\textwidth}{|X|X|X|}
  \hline
  $38 \times 47$ & $84 \times 73$ & $66 \times 52$ \\
  \rule[-3.3em]{0pt}{4.1em} & & \\
  \hline
\end{tabularx}

% -- E6 ---------------------------------------------------------------------
\qhead[cbgreen]{E6}{Three digits by two digits}

\begin{remember}[Nothing new here]
  Three digits by two is still \key{two rows}. Row 1 is the whole number times the
  units. Row 2 is the whole number times the tens. It is only longer, not harder.
\end{remember}

\vspace{2pt}
\noindent
\begin{tabularx}{\textwidth}{|X|X|X|}
  \hline
  $124 \times 32$ & $236 \times 18$ & $407 \times 25$ \\
  \rule[-3.7em]{0pt}{4.5em} & & \\
  \hline
\end{tabularx}

\vspace{9pt}
\noindent
\begin{tabularx}{\textwidth}{|X|X|X|}
  \hline
  $318 \times 46$ & $529 \times 37$ & $324 \times 18$ \\
  \rule[-3.7em]{0pt}{4.5em} & & \\
  \hline
\end{tabularx}

\vspace{9pt}
\noindent\textbf{Look back.} The last one, $324 \times 18$, is the second half of a
cube root: $18 \times 18 = 324$, and then $324 \times 18$. So
$\sqrt[3]{\blank[2.2cm]} = 18$.

% ===========================================================================
% PART 3 -- LONG DIVISION
% ===========================================================================
\hwsection[cbgreen]{Part 3 \textperiodcentered\ Long division: bring down the next digit}

\begin{wordhelp}[Word help --- the words in a division]
  {\renewcommand{\arraystretch}{1.15}%
  \begin{tabularx}{\linewidth}{@{}l@{\hspace{1.4em}}X@{}}
    \eng{divisor}     & the number you are dividing \textbf{by}. In $966 \div 23$ it is 23. \\
    \eng{goes into}   & ``23 goes into 96 four times'' means $96 \div 23 = 4$, with some left over. \\
    \eng{bring down}  & take the next digit of the big number and write it beside what is left. \\
    \eng{remainder}   & what is left at the very end. We write $75 \div 8 = 9$ r 3. \\
  \end{tabularx}}
\end{wordhelp}

\begin{remember}[Divide, multiply, subtract, bring down --- then start again]
  \noindent
  \begin{minipage}[t]{0.30\linewidth}
    \vspace{4pt}
    {\renewcommand{\arraystretch}{1.15}%
    \begin{tabular}{r@{}c@{\;}c@{\;}c}
                        &   & 4 & 2 \\ \cline{2-4}
      $23\,\big)$       & 9 & 6 & 6 \\
                        & 9 & 2 &   \\ \cline{2-3}
                        &   & 4 & 6 \\
                        &   & 4 & 6 \\ \cline{3-4}
                        &   &   & 0 \\
    \end{tabular}}
  \end{minipage}%
  \begin{minipage}[t]{0.70\linewidth}
    \vspace{2pt}
    \key{Divide:} 23 goes into 96 \textbf{four} times. Write the 4 on top, above the 6. \\
    \key{Multiply:} $4 \times 23 = 92$. Write it underneath. \\
    \key{Subtract:} $96 - 92 = 4$. \\
    \key{Bring down} the next digit, the 6, to make \textbf{46}. \\[3pt]
    Now start again: 23 goes into 46 \textbf{twice}, $2 \times 23 = 46$,
    $46 - 46 = 0$. Nothing is left, so $966 \div 23 = \mathbf{42}$.
  \end{minipage}
\end{remember}

% -- E7 ---------------------------------------------------------------------
\qhead[cbgreen]{E7}{Warm up: one-digit divisors, no remainder}

\vspace{2pt}
\begin{multicols}{2}
\begin{enumerate}[label=(\alph*), itemsep=12pt]
  \item $252 \div 6 = \blank[2.2cm]$
  \item $469 \div 7 = \blank[2.2cm]$
  \item $584 \div 8 = \blank[2.2cm]$
  \item $765 \div 9 = \blank[2.2cm]$
  \item $1344 \div 4 = \blank[2.2cm]$
  \item $2016 \div 6 = \blank[2.2cm]$
  \item $3528 \div 7 = \blank[2.2cm]$
  \item $5832 \div 8 = \blank[2.2cm]$
\end{enumerate}
\end{multicols}

% -- E8 ---------------------------------------------------------------------
\qhead[cbgreen]{E8}{One-digit divisors, with a remainder}

\noindent Write each answer as a whole number and a remainder, like this:
\quad $75 \div 8 = 9$ r 3.

\vspace{5pt}
\begin{multicols}{2}
\begin{enumerate}[label=(\alph*), itemsep=12pt]
  \item $437 \div 6 = \blank[2.8cm]$
  \item $529 \div 7 = \blank[2.8cm]$
  \item $815 \div 9 = \blank[2.8cm]$
  \item $1000 \div 7 = \blank[2.8cm]$
  \item $2345 \div 8 = \blank[2.8cm]$
  \item $4567 \div 9 = \blank[2.8cm]$
\end{enumerate}
\end{multicols}

% -- E9 ---------------------------------------------------------------------
\qhead[cbgreen]{E9}{Two-digit divisors --- set it out properly}

% The one-line instruction goes BEFORE the box, not after it. This heading lands at
% the foot of a page and one line fits there; the box does not, so leading with the
% box left the heading stranded on its own.
\noindent These all come out exactly. Show the full method in each box.

\begin{remember}[Write the list first]
  Before you start, write the first six multiples of the divisor in the corner of
  the box. Then every ``how many times does it go in'' is \key{reading}, not
  guessing.
\end{remember}

\vspace{4pt}
\noindent
\begin{tabularx}{\textwidth}{|X|X|}
  \hline
  $768 \div 16$ & $851 \div 23$ \\
  \rule[-4.0em]{0pt}{4.8em} & \\
  \hline
\end{tabularx}

\vspace{9pt}
\noindent
\begin{tabularx}{\textwidth}{|X|X|}
  \hline
  $1224 \div 24$ & $2072 \div 37$ \\
  \rule[-4.0em]{0pt}{4.8em} & \\
  \hline
\end{tabularx}

\vspace{9pt}
\noindent
\begin{tabularx}{\textwidth}{|X|X|}
  \hline
  $4212 \div 18$ & $3375 \div 15$ \\
  \rule[-4.0em]{0pt}{4.8em} & \\
  \hline
\end{tabularx}

\vspace{9pt}
\noindent
\begin{tabularx}{\textwidth}{|X|X|}
  \hline
  $1088 \div 32$ & $1222 \div 47$ \\
  \rule[-4.0em]{0pt}{4.8em} & \\
  \hline
\end{tabularx}

\vspace{9pt}
\noindent
\begin{tabularx}{\textwidth}{|X|X|}
  \hline
  $986 \div 29$ & $2408 \div 43$ \\
  \rule[-4.0em]{0pt}{4.8em} & \\
  \hline
\end{tabularx}

\vspace{9pt}
\noindent
\begin{tabularx}{\textwidth}{|X|X|}
  \hline
  $1872 \div 26$ & $2940 \div 35$ \\
  \rule[-4.0em]{0pt}{4.8em} & \\
  \hline
\end{tabularx}

% -- E10 --------------------------------------------------------------------
\qhead[cbgreen]{E10}{Two-digit divisors, with a remainder}

% A one-line instruction rather than a remember box. The box made this question a
% 10cm atom that would not start at the foot of a page, and stranded a third of
% one every time.
\noindent The remainder is always \key{smaller than the divisor}. Left with 30 after
dividing by 24? Then 24 goes in \key{one more time} --- go back.

\vspace{6pt}
\noindent
\begin{tabularx}{\textwidth}{|X|X|}
  \hline
  $1000 \div 24$ & $875 \div 32$ \\
  \rule[-3.4em]{0pt}{4.2em} & \\
  \hline
\end{tabularx}

\vspace{9pt}
\noindent
\begin{tabularx}{\textwidth}{|X|X|}
  \hline
  $2500 \div 36$ & $1500 \div 45$ \\
  \rule[-3.6em]{0pt}{4.4em} & \\
  \hline
\end{tabularx}

\vspace{9pt}
\noindent
\begin{tabularx}{\textwidth}{|X|X|}
  \hline
  $3000 \div 28$ & $1800 \div 26$ \\
  \rule[-3.6em]{0pt}{4.4em} & \\
  \hline
\end{tabularx}

% -- E11 --------------------------------------------------------------------
\qhead[cbgreen]{E11}{Find the missing number}

\begin{remember}[Every multiplication is two divisions]
  If $24 \times 51 = 1224$, then $1224 \div 24 = 51$ \key{and} $1224 \div 51 = 24$.
  Look at \key{which} number is missing before you decide what to do: if the
  \textbf{big} number is missing, you \key{multiply}.
\end{remember}

\vspace{2pt}
\begin{multicols}{2}
\begin{enumerate}[label=(\alph*), itemsep=12pt]
  \item $24 \times \blank[1.8cm] = 1224$
  \item $\blank[1.8cm] \times 16 = 768$
  \item $2072 \div \blank[1.8cm] = 56$
  \item $\blank[1.8cm] \div 18 = 234$
  \item $37 \times \blank[1.8cm] = 851$
  \item $\blank[1.8cm] \times 18 = 324$
  \item $3375 \div \blank[1.8cm] = 225$
  \item $\blank[1.8cm] \div 8 = 729$
  \item $43 \times \blank[1.8cm] = 2408$
  \item $\blank[1.8cm] \times 26 = 1872$
  \item $2940 \div \blank[1.8cm] = 84$
  \item $\blank[1.8cm] \div 29 = 34$
  \item $47 \times \blank[1.8cm] = 1222$
  \item $\blank[1.8cm] \times 32 = 1088$
  \item $4212 \div \blank[1.8cm] = 234$
  \item $\blank[1.8cm] \div 15 = 225$
\end{enumerate}
\end{multicols}

% -- E12 --------------------------------------------------------------------
% No boxes on this one on purpose. Everything here can be done in the head or in
% the margin, and the point is speed and choosing the method, not setting it out.
\qhead[cbgreen]{E12}{Mixed drill --- no boxes, no help}

\noindent Both operations, mixed up. Work down the columns as fast as you can be
accurate. \key{Two pairs of these have the same answer --- find both pairs.}

% A tabular, not multicols. Twenty-four items in three columns will not fit in the
% tail of a page, and multicols splits itself across the break and then refuses to
% let anything follow it -- which stranded a whole page. A tabular moves whole.
\vspace{8pt}
\noindent
\begin{tabular}{@{}p{0.32\textwidth}p{0.32\textwidth}p{0.30\textwidth}@{}}
  1.\ $480 \div 6 = \blank[1.4cm]$   & 11.\ $936 \div 12 = \blank[1.4cm]$  & 21.\ $999 \div 3 = \blank[1.4cm]$ \\[12pt]
  2.\ $24 \times 30 = \blank[1.4cm]$ & 12.\ $18 \times 15 = \blank[1.4cm]$ & 22.\ $44 \times 25 = \blank[1.4cm]$ \\[12pt]
  3.\ $3600 \div 9 = \blank[1.4cm]$  & 13.\ $2400 \div 16 = \blank[1.3cm]$ & 23.\ $864 \div 12 = \blank[1.4cm]$ \\[12pt]
  4.\ $45 \times 12 = \blank[1.4cm]$ & 14.\ $27 \times 40 = \blank[1.4cm]$ & 24.\ $125 \times 8 = \blank[1.4cm]$ \\[12pt]
  5.\ $720 \div 8 = \blank[1.4cm]$   & 15.\ $546 \div 7 = \blank[1.4cm]$   & 25.\ $756 \div 7 = \blank[1.4cm]$ \\[12pt]
  6.\ $16 \times 25 = \blank[1.4cm]$ & 16.\ $23 \times 30 = \blank[1.4cm]$ & 26.\ $48 \times 25 = \blank[1.4cm]$ \\[12pt]
  7.\ $1000 \div 8 = \blank[1.4cm]$  & 17.\ $1250 \div 25 = \blank[1.3cm]$ & 27.\ $630 \div 15 = \blank[1.4cm]$ \\[12pt]
  8.\ $60 \times 14 = \blank[1.4cm]$ & 18.\ $19 \times 21 = \blank[1.4cm]$ & 28.\ $26 \times 40 = \blank[1.4cm]$ \\[12pt]
  9.\ $812 \div 4 = \blank[1.4cm]$   & 19.\ $4000 \div 25 = \blank[1.3cm]$ & 29.\ $1440 \div 12 = \blank[1.3cm]$ \\[12pt]
  10.\ $35 \times 11 = \blank[1.4cm]$ & 20.\ $32 \times 15 = \blank[1.4cm]$ & 30.\ $55 \times 12 = \blank[1.4cm]$ \\
\end{tabular}

% ===========================================================================
% PART 4 -- USE IT
% ===========================================================================
\hwsection[cbgreen]{Part 4 \textperiodcentered\ Use it}

\begin{wordhelp}[Word help --- which operation does the sentence want?]
  \eng{each} / \eng{every} --- usually multiply. 32 boxes with 124 in \textbf{each} is
  $32 \times 124$. \\
  \eng{in rows of} \sep \eng{shared equally between} --- divide.
  \quad \eng{how many are left over} --- the remainder.
\end{wordhelp}

\noindent Write the calculation first, then the answer.

\begin{enumerate}[label=\textbf{\arabic*.}, itemsep=11pt]

  \item \textbf{The hall.} There are \textbf{1224 chairs} in the hall, set out in
        rows of \textbf{24}. How many rows are there?
        \begin{workspace}[3.0cm]\end{workspace}

  \item \textbf{The books.} Mr Bowen orders \textbf{32 boxes} of books with
        \textbf{124 books} in each box. How many books is that altogether?
        \begin{workspace}[3.0cm]\end{workspace}

  \item \textbf{The pencils.} Mr Bowen has \textbf{1000 pencils} and shares them
        equally between \textbf{24 classes}. How many does each class get, and how
        many are left over?
        \begin{workspace}[3.2cm]\end{workspace}

  \item \textbf{The order.} The school orders \textbf{18 packs} of pens with
        \textbf{236 pens} in each pack. The pens are then shared equally between
        \textbf{24 classrooms}. How many pens does each classroom get?
        \begin{workspace}[3.6cm]\end{workspace}

  % The round-up question. 61 r 2 boxes is not an answer -- you cannot buy part of
  % a box -- and this is the only place on the sheet where the remainder makes the
  % answer go UP.
  \item \textbf{The tiles.} Mr Bowen tiles a floor with \textbf{47 rows} of
        \textbf{26 tiles}. The tiles come in boxes of \textbf{20}. How many tiles
        does he need, and how many boxes must he buy?
        \begin{workspace}[3.6cm]\end{workspace}

  % The deadpan one, and the last thing on the sheet: two steps, and the leftover
  % space at the foot of the page belongs to the student, not to the margin.
  \item \textbf{The stickers.} Mr Bowen has \textbf{4212 stickers} and sticks
        \textbf{18} of them on every page of a sticker book. Each book has
        \textbf{50 pages}. How many pages does he fill, and how many \textbf{whole
        books} is that?
        \begin{workspace}[4.4cm]\end{workspace}

\end{enumerate}

% ===========================================================================
% ANSWER KEY (teacher copy only)
% ===========================================================================
\ifkey
\newpage
\hwsection[cbred]{Answer Key --- Teacher Copy}

\linespread{1.0}\selectfont
\setlist[enumerate]{itemsep=2pt, topsep=3pt, leftmargin=*}
\setlist[itemize]{itemsep=2pt, topsep=3pt, leftmargin=*}

\qhead[cbred]{E1}{Fill in the missing number}
\noindent (a) 42 \quad (b) 56 \quad (c) 54 \quad (d) 48 \quad (e) 72 \quad (f) 49
\quad (g) 36 \quad (h) 64 \\
(i) 63 \quad (j) 36 \quad (k) 81 \quad (l) 32 \quad (m) 84 \quad (n) 88 \quad
(o) 108 \quad (p) 28

\smallskip
\noindent{\small Time this one rather than marking it. Under two minutes for all
sixteen and the tables are not what is slowing them down --- look at Part 3 instead.
The three worth re-drilling are (b) $8 \times 7$, (e) $9 \times 8$ and (i)
$9 \times 7$: they turn up in both rows of a long multiplication \emph{and} in every
long division by 7, 8 or 9, so an error there is paid for three times.}

\qhead[cbred]{E2}{Multiply by a ten}
\noindent (a) 260 \quad (b) 780 \quad (c) 940 \quad (d) 2080 \quad (e) 1080
\quad (f) 3650

\smallskip
\noindent{\small This is the second row of E5 taken out and practised alone, which is
why it comes before the method. The error to look for is a \textbf{missing zero} ---
78 instead of 780 in (b). A student who makes it here makes it in every question in
E5, and it is far faster to catch on this line than inside a four-line sum.}

\qhead[cbred]{E3}{Count up in a two-digit number}
\noindent
\begin{tabularx}{\linewidth}{@{}p{1.2cm}X@{}}
  $23$ & 23, 46, 69, 92, 115, 138 \\
  $16$ & 16, 32, 48, 64, 80, 96 \\
  $24$ & 24, 48, 72, 96, 120, 144 \\
  $37$ & 37, 74, 111, 148, 185, 222 \\
\end{tabularx}

\smallskip
\noindent{\small This question is the whole sheet in miniature and it is the one to
insist on. A student who writes the list before dividing is \emph{reading} at every
step of a long division; a student who does not is guessing, rubbing out and
guessing again. Note that these four are exactly the divisors in E9 --- 23 and 16 and
24 --- so a student who does this table properly has already done half of E9.}

\qhead[cbred]{E4}{Warm up: two digits by one digit}
\noindent $53 \times 7 = 371$ \quad $68 \times 6 = 408$ \quad $87 \times 9 = 783$
\quad $76 \times 8 = 608$ \\
$94 \times 4 = 376$ \quad $47 \times 8 = 376$ \quad $65 \times 7 = 455$ \quad
$89 \times 6 = 534$ \\
$58 \times 9 = 522$ \quad $73 \times 6 = 438$ \quad $96 \times 7 = 672$ \quad
$84 \times 5 = 420$

\smallskip
\noindent{\small The error here is carrying, not the tables. $87 \times 9$ carries at
both columns and is the one to check. \\
$94 \times 4$ and $47 \times 8$ both come to \textbf{376}, and that is not a
coincidence: halving one number and doubling the other leaves the answer alone. A
student who spots it has found a genuinely useful shortcut --- $16 \times 25$ in E12
is the same trick, and becomes $8 \times 50$, then $4 \times 100$.}

\qhead[cbred]{E5}{Two digits by two digits}
\noindent
\begin{tabularx}{\linewidth}{@{}p{4.4cm}p{4.4cm}X@{}}
  $32 \times 14 = 448$  & $45 \times 23 = 1035$ & $28 \times 36 = 1008$ \\
  $57 \times 42 = 2394$ & $63 \times 27 = 1701$ & $86 \times 35 = 3010$ \\
  $74 \times 58 = 4292$ & $49 \times 61 = 2989$ & $95 \times 24 = 2280$ \\
  $38 \times 47 = 1786$ & $84 \times 73 = 6132$ & $66 \times 52 = 3432$ \\
\end{tabularx}

\smallskip
\noindent{\small Mark the \emph{working}, not just the answer. \textbf{One row of
working is the whole diagnosis}: the student multiplied by the units digit and
stopped, so the answer is roughly a tenth of the right size --- 448 becomes 128. That
is a method error and needs re-teaching. An answer out by a small amount with two
rows present is arithmetic, and needs a re-check, not a re-teach.}

\qhead[cbred]{E6}{Three digits by two digits}
\noindent $124 \times 32 = 3968$ \quad $236 \times 18 = 4248$ \quad
$407 \times 25 = 10175$ \\
$318 \times 46 = 14628$ \quad $529 \times 37 = 19573$ \quad
$324 \times 18 = 5832$

\smallskip
\noindent \textbf{Look back:} $\sqrt[3]{5832} = 18$.

\smallskip
\noindent{\small $407 \times 25$ is the one with the \textbf{zero in the middle} and
it is the interesting failure: students skip the zero column and shorten the row by a
digit. \\
$324 \times 18$ is not decoration. $18 \times 18 = 324$, so this multiplication
\emph{is} the second half of proving a cube root by hand, and a student who reaches
5832 has done something worth saying out loud. It also reappears as E7(h),
$5832 \div 8$, from the other direction.}

\qhead[cbred]{E7}{One-digit divisors, no remainder}
\noindent (a) 42 \quad (b) 67 \quad (c) 73 \quad (d) 85 \quad (e) 336 \quad
(f) 336 \quad (g) 504 \quad (h) 729

\smallskip
\noindent{\small (e) and (f) both come to 336 --- if a student notices, that is a good
sign and worth a word. (h) $5832 \div 8 = 729$ is the same number as E6's last
multiplication seen from the other end, and 729 is $9^3$; a class that spots that has
understood something about Unit 1 that no exercise can be set for.}

\qhead[cbred]{E8}{One-digit divisors, with a remainder}
\noindent (a) 72 r 5 \quad (b) 75 r 4 \quad (c) 90 r 5 \quad (d) 142 r 6 \quad
(e) 293 r 1 \quad (f) 507 r 4

\smallskip
\noindent{\small The diagnosis is in the remainder. A remainder \textbf{bigger than
the divisor} is a method error --- the student stopped dividing too early --- and
needs re-teaching. A remainder wrong by one or two is arithmetic and needs a
re-check. Only the first is worth class time. \\
(c) and (f) are the two with a \textbf{zero in the answer}: 815 divided by 9 gives 9,
then nothing, then 0 --- and the zero must be written. A student who answers ``9 r 5''
has dropped a digit, and that is the single most common long-division error there is.}

\qhead[cbred]{E9}{Two-digit divisors --- set it out properly}
\noindent $768 \div 16 = \mathbf{48}$ \quad $851 \div 23 = \mathbf{37}$ \quad
$1224 \div 24 = \mathbf{51}$ \\
$2072 \div 37 = \mathbf{56}$ \quad $4212 \div 18 = \mathbf{234}$ \\
$3375 \div 15 = \mathbf{225}$ \quad $1088 \div 32 = \mathbf{34}$ \\
$1222 \div 47 = \mathbf{26}$ \quad $986 \div 29 = \mathbf{34}$ \\
$2408 \div 43 = \mathbf{56}$ \quad $1872 \div 26 = \mathbf{72}$ \quad
$2940 \div 35 = \mathbf{84}$

\smallskip
\noindent{\small Look in the corner of the box before you look at the answer. A
student who wrote the multiples list first and got it wrong has made an arithmetic
slip; a student with no list who got it wrong has been guessing, and telling them the
answer teaches them nothing. \\
$1224 \div 24$ has a \textbf{zero step} in the middle of the method --- 24 does not go
into 12, so you take 122 --- and it is the one that stalls people. $3375 \div 15$ is
the hardest, and 225 is $15^2$, which is the patio from the lesson.}

\qhead[cbred]{E10}{Two-digit divisors, with a remainder}
\noindent $1000 \div 24 = \mathbf{41}$ r $\mathbf{16}$ \quad
$875 \div 32 = \mathbf{27}$ r $\mathbf{11}$ \quad
$2500 \div 36 = \mathbf{69}$ r $\mathbf{16}$ \\
$1500 \div 45 = \mathbf{33}$ r $\mathbf{15}$ \quad
$3000 \div 28 = \mathbf{107}$ r $\mathbf{4}$ \quad
$1800 \div 26 = \mathbf{69}$ r $\mathbf{6}$

\smallskip
\noindent{\small Check the remainders first: any remainder of 24 or more in the first
one is a stopped-too-early error, and the student can find it themselves if you ask
``can 24 go in again?''. That question is worth more than the correction. \\
$3000 \div 28$ is the hard one, and it is hard for a reason worth naming: 28 does not
go into 30 twice, it goes once, and the second digit of the answer is a \textbf{0}.
Expect ``17 r 4''. \\
$1000 \div 24$ is the same division as Part 4 question 3, so a student who has done it
here has done it there.}

\qhead[cbred]{E11}{Find the missing number}
\noindent (a) 51 \quad (b) 48 \quad (c) 37 \quad (d) 4212 \quad (e) 23 \quad
(f) 18 \quad (g) 15 \quad (h) 5832 \\
(i) 56 \quad (j) 72 \quad (k) 35 \quad (l) 986 \quad (m) 26 \quad (n) 34 \quad
(o) 18 \quad (p) 3375

\smallskip
\noindent{\small (d), (h), (l) and (p) are the hard ones, because the missing number is
the \emph{big} one and the student has to \textbf{multiply} rather than divide:
$234 \times 18$, $729 \times 8$, $34 \times 29$, $225 \times 15$. A student who tries
to divide there has not looked at which number is missing, and that is a reading error,
not an arithmetic one. \\
Every number in this question appears somewhere else on the sheet, so it is also a
check on whether the earlier work was done or copied --- a student who did E9 properly
can fill in half of this from memory.}

\qhead[cbred]{E12}{Mixed drill --- no boxes, no help}
\noindent 1.\ 80 \quad 2.\ 720 \quad 3.\ 400 \quad 4.\ 540 \quad 5.\ 90 \quad
6.\ 400 \quad 7.\ 125 \quad 8.\ 840 \quad 9.\ 203 \quad 10.\ 385 \\
11.\ 78 \quad 12.\ 270 \quad 13.\ 150 \quad 14.\ 1080 \quad 15.\ 78 \quad
16.\ 690 \quad 17.\ 50 \quad 18.\ 399 \quad 19.\ 160 \quad 20.\ 480 \\
21.\ 333 \quad 22.\ 1100 \quad 23.\ 72 \quad 24.\ 1000 \quad 25.\ 108 \quad
26.\ 1200 \quad 27.\ 42 \quad 28.\ 1040 \quad 29.\ 120 \quad 30.\ 660

\smallskip
\noindent \textbf{The two pairs:} 3 and 6 both make \textbf{400}; 11 and 15 both make
\textbf{78}.

\smallskip
\noindent{\small Time this one and mark it as a block --- it is a speed question, not
a method question, and the whole diagnosis is \emph{where they slowed down}. A student
who is fast on the multiplications and slow on the divisions has the tables but not
the division method, which is the opposite way round from what most people assume. \\
9 ($812 \div 4 = 203$) and 25 ($756 \div 7 = 108$) are the \textbf{zero in the answer}
again, and 7 ($1000 \div 8 = 125$) is the one they will want a calculator for. The
pairs are there to make them look back over their own work, which is the check they
never do voluntarily.}

\qhead[cbred]{E13}{Use it}
\begin{enumerate}[label=\arabic*., itemsep=2pt]
  \item $1224 \div 24 = \mathbf{51}$ rows.
  \item $32 \times 124 = \mathbf{3968}$ books.
  \item $1000 \div 24 = 41$ r 16: \textbf{41 pencils} each, with \textbf{16 left
        over}.
  \item $18 \times 236 = 4248$ pens, then $4248 \div 24 = \mathbf{177}$ pens each.
  \item $47 \times 26 = \mathbf{1222}$ tiles. $1222 \div 20 = 61$ r 2, so he must buy
        \textbf{62 boxes}.
  \item $4212 \div 18 = 234$ pages. Then $234 \div 50 = 4$ r 34, so \textbf{234
        pages}, which is \textbf{4 whole books} (with 34 pages of a fifth).
\end{enumerate}
\noindent{\small 1 and 2 are E9 and E6 wearing a story, so a student who can do those
and not these has a reading gap, not an arithmetic one --- read the sentence back to
them rather than reworking the sum. \\
3 is the sentence test: ``41.67'' is a correct division answering a question nobody
asked, because you cannot give a class two-thirds of a pencil. \\
4 is $236 \times 18$ from E6 followed by a division, and it comes out \emph{exactly},
which is the reward for getting the multiplication right. A student who reaches 4248
and then a messy remainder has an error in step one, not step two --- send them back
to the multiplication rather than to the division. \\
\textbf{5 is the one to spend time on.} $47 \times 26 = 1222$ is E9 seen from the
other side, but the second half is the only place on the sheet where the remainder
makes the answer go \textbf{up}: 61 boxes is not enough, because the last 2 tiles
still have to come in a box. Expect ``61'' and ``61 r 2'' from most of the class, and
treat it as a discussion rather than a cross --- the arithmetic is right and the
reading is not. \\
6 is the last thing on the sheet and the only one where the remainder has to be
\emph{read} rather than reported: 4 r 34 means four full books and part of a fifth,
not ``4.68 books''. The common failure is stopping at 234 and calling those books.}

\fi
\end{document}
